% 全期間累積の変更履歴（課題3 用）
% revision_history/revision_history.tex から \input される．
%
% ライフサイクル:
%   - 提出のタイミングで，sections/changelog_current.tex の履歴行を
%     このファイルの末尾（\endhead と \bottomrule の間）にコピーする
%   - このファイルは「リセットしない」．全期間の履歴がここに蓄積される
%
% 書式:
%   YYYY-MM-DD & 学籍番号 : 氏名 & 対象（章・節） & 変更内容 \\
%
% 前提パッケージ（呼び出し側でロード済みであること）:
%   \usepackage{longtable}
%   \usepackage{booktabs}
%   \usepackage{array}

\begin{longtable}{p{2.2cm}p{3.5cm}p{4cm}p{5.5cm}}
\toprule
日付 & 担当 & 対象（章・節） & 変更内容 \\
\midrule
\endfirsthead

\toprule
日付 & 担当 & 対象（章・節） & 変更内容 \\
\midrule
\endhead

% --- 全期間の履歴行（提出時に changelog_current.tex から追記される） ---
2026-05-29 & 25G1099 : 長山 魁良 & 第3章 担当C詳細設計（オブジェクト設計・関数設計） & シリアル通信内容に音の長さを追加し，変数一覧・関数一覧を更新 \\
2026-05-29 & 24G1056 : 慶田 優 & 第1章 作業計画（WBS） & 担当AにProcessing GUI操作・シリアル通信設計を追加，担当Bに予備振り前のキャリブレーション状態・輝度エンコード制御の設計を追加 \\
2026-05-29 & 24G1056 : 慶田 優 & 第1章 作業計画（ガントチャート） & 機材購入遅延により，ハードウェア先行からソフトウェア設計先行（→ハードウェア→テスト）のフローへ全体を変更 \\
2026-05-30 & 25G1098 : 長野 志哉 & 第2章 担当B基本設計（操作インターフェース設計，状態遷移図） & フォトトランジスタのアナログ読み取りへの変更，CALIBRATE状態の追加に伴い内容を修正． \\
2026-05-30 & 25G1098 : 長野 志哉 & 第3章 担当B詳細設計 （ハードウェア設計，主要変数，担当 C との共有変数，主要関数，処理フローの設計）& 全体的な変数および関数名の変更．フォトトランジスタのアナログ読み取りへの変更に伴い拍検知の方法を変更し，detectPhotoBeat関数を追加．CALIBRATE状態の追加に伴いupdateState関数を修正．フローチャートを重要な関数のみに厳選． \\
2026-06-12 & 25G1098 : 長野 志哉 & 第2章 担当B基本設計 （テンポ推定）& 演奏開始制御方式を拍カウントによるオフセット制御からフォトトランジスタによる明るさを基準にした判定に変更． \\
2026-06-12 & 25G1098 : 長野 志哉 & 第3章 担当B詳細設計 （ハードウェア設計）& 演奏開始制御方式を拍カウントによるオフセット制御からフォトトランジスタによる明るさを基準にした判定に変更． \\
2526-06-12 & 25G1033 : 長見 東磨 & 第2章 担当A基本設計（システム構成図・操作インターフェース設計・状態遷移図） & 楽器番号制御仕様に画像差し替え \\
2026-06-12 & 25G1033 : 長見 東磨 & 第3章 担当A詳細設計（回路設計・処理フロー設計） & 楽器番号制御仕様に画像差し替え \\
2026-06-12 & 25G1033 : 長見 東磨 & 第3章 担当A詳細設計（シリアル通信コマンド仕様） & BPM下限・上限値を「60〜200」から「60〜120」へ修正 \\ 
2026-06-19 & 25G1099 : 長山 魁良 & 第2章 担当C基本設計（シリアル通信パケット仕様） & 音の長さ(LONG)の範囲を「0〜127」から「2，4，8」に変更\\
2026-06-19 & 24G1056 : 慶田 優 & 第1章 評価指標（TPM一覧・TPM監視体制） & テンポ予測誤差をTPM-6として追加（目標値$\leq 1\,\mathrm{ms}$，監視方法：100回計測の平均誤差記録） \\
2026-06-19 & 24G1056 : 慶田 優 & 第6章 テスト実施記録 & 担当B/C/D結合テスト（INT-2）の結果を追記（カエルの歌演奏確認・合格） \\
2026-06-22 & 25G1033 : 長見 東磨 & 第1章 担当A（部品一覧） & サーボモータをSG92RからMG996Rに変更，LCD項目を削除 \\
2026-06-22 & 25G1033 : 長見 東磨 & 第3章 担当A（ピン接続・回路設計） & LCD無し仕様に修正・画像差し替え \\
2026-06-22 & 25G1033 : 長見 東磨 & 第3章 担当A（関数設計） & 変数・関数一覧を更新 \\


\bottomrule
\endfoot
\end{longtable}
